\section{Man's Place in Nature}

In 1863, \textbf{T. H. Huxley} published \emph{Man's Place in Nature} to demonstrate that man is merely an animal lacking any higher or transcendent quality.

\paragraph{Microcosm and Macrocosm}
\textbf{Valentin Tomberg} explains that ``initiation" properly means ``knowledge of the beginning". This implies that there are actually two initiations:

\begin{itemize}
\item \textbf{Hermetic Initiation} is the conscious experience of the initial microcosmic state. 
\item \textbf{Pythagorean Initiation} is the conscious experience of the initial macrocosmic state 
\end{itemize}
Hermetic initiation is achieved when the initiate descends into the depths of the human being, i.e., becomes aware of himself as the image and likeness of God, as the Intellect. This requires an alchemical transformation, i.e., a real change of substance.

Pythagorean initiation is achieved by a ``rapture", a going out of oneself (rather than within), so that the macrocosmic, or heavenly, layers reveal themselves. The initiate becomes conscious of the musical and mathematical structure of the macrocosm.

\paragraph{The Natural Law of Life}
In order to get to the beginning, the initiate must start from his present position. Therefore, the initiate studies the inner structure of consciousness in order to penetrate to the beginning (i.e., Eden). We also need to study the structure of the universe and man's place in it, in order to transcend it (i.e., Paradise). Unlike profane cosmology which studies the quantitative properties of the objects in relation to each other, esoteric cosmology describes the universe qualitatively, as it relates to consciousness.

This is what \textbf{Boris Mouravieff} tries to do, by drawing on the resources of the Western Tradition. Because of the dual situation, man's existence has two primary purposes:

\begin{itemize}
\item As part of the whole, he serves the aims of the cosmos. 
\item As an isolated individual, he can pursue his own aims. 
\end{itemize}
The General Law refers to the cosmic forces that tend to keep man in his place in nature, that is, as part of organic life on earth, the same as the plants and animals. This is what makes Huxley's thesis plausible, even ``obvious" to educated men of our day. This is the result of the forgetting of Being, i.e., that the natural world is one link in a Chain of Being.

Examples of the General Law:

\begin{itemize}
\item \textbf{Hunger}: the toil necessary to assure our subsistence, not to mention the satisfaction of other desires. 
\item \textbf{Sex}: sexual instinct and its consequences in procreation and parenthood. 
\item \textbf{Fear}: fear of death, injury, destitution, etc. 
\end{itemize}
Such motivations take up most, if not all, of man's time and energy. A dispassionate and objective search for truth, the very definition of Intelligence, becomes nearly impossible. On the contrary, the definition of intelligence becomes inverted to mean that the man who is most successful in addressing these aspects of the General Law is regarded as the most intelligent.

Hence, the goals of achieving financial security, raising a family, becoming healthy, and the like, are regarded as the most important. Even religious sensibilities and political opinions are oriented primarily to those goals. I don't think I need to provide specific examples as they are so commonplace; on the contrary, it would be difficult to provide examples from politics and religion of achieving higher goals.

In the scheme of nature, life becomes an incessant series of minor battles:

\begin{itemize}
\item Seeking the pleasant while avoiding the unpleasant 
\item Loving what feels good and hating what feels bad 
\item Desire followed by delight at its satisfaction or frustration 
\item Fear of the worst while hoping for the best 
\end{itemize}
The life of bourgeois happiness is summed up in winning those battles. Nevertheless, there are still occasions for heroic actions, since life presents so many challenges.

\begin{itemize}
\item It takes courage to face up to and overcome difficult situations 
\item The hope of success overcomes the fear of failure 
\end{itemize}
Sometimes boundary situations may be so overwhelming, that a man or woman becomes aware of something higher. For them, exoteric religion may be the vehicle for transcendence.

Now this situation is entirely lawful, since if everyone escaped from the General Law, it would mean the end of the world. That is why esoterism is only for the few who are both capable of understanding its principles and have the opportunity to pursue it.

\paragraph{Macrocosmic Initiation}
Macrocosmic initiation is the realization that the natural life is not the whole story. The human being is not only natural but is also a person, transcending nature. The notion that the natural world is all that exists is considered ``realistic". Success is measured in the satisfaction of desires and the reduction in the causes of fear. This is considered ``life affirming" (although not necessarily in the Nietzschean sense) in opposition to the ``life denying" attitude of ascetical practices to minimize the impact of the forces of the General Law.

If survival is the goal of the natural life, then salvation and liberation are the goals of a higher life. \textbf{Frithjof Schuon}, expanding on \textbf{Rene Guenon}, gives us the examples of the Noble Man and the Holy Man as the fulfillment of the latter two goals.

\subparagraph{The Noble Man}
\begin{quotex}
The noble man is one who dominates himself. The noble man is one who masters himself and loves to master himself; the base man is one who does not master himself and shrinks in horror from mastering himself. The noble man always maintains himself at the centre; he never loses sight of the symbol, the spiritual gift of things, the sign of God, a gratitude that is both ascending and radiating. The noble man is naturally detached from mean things, sometimes against his own interests; and he is naturally generous through greatness of soul. [\emph{Esoterism as Principle and as Way}] 

\end{quotex}
\subparagraph{The Holy Man}
\begin{quotex}
Transcending oneself: this is the great imperative of the human condition; and there is another that anticipates it and at the same time prolongs it: dominating oneself. The noble man is one who dominates himself; the holy man is one who transcends himself. \emph{Nobility and holiness are the imperatives of the human state}.

Intelligence, since it distinguishes, has the faculty of perceiving proportions. The spiritual man integrates these proportions into his will, into his soul and into his life. All defects manifest a lack of proportion; they are errors that are lived.

To be spiritual means not to deny with one's `being' what one affirms with one's `knowledge', that is to say, what is accepted by the intelligence. Truth lived: incorruptibility and generosity. [\emph{Islam and The Destiny of Man}] 

\end{quotex}


\flrightit{Posted on 2014-10-22 by Cologero }

\begin{center}* * *\end{center}

\begin{footnotesize}\begin{sffamily}



\texttt{Tom Blanchard on 2014-10-23 at 00:00 said: }

With regard to Pythagorean initiation, a good primer on the subject is HRH The Prince of Wales' book, ``Harmony" ( \url{http://www.amazon.com/Harmony-New-Way-Looking-World/dp/0007348037/} ), which outlines the origins and nature of sacred geometry and music in a way that is very approachable for the uninitiated reader. Prince Charles is a follower of Guenon ( \url{http://www.sacredweb.com/conference06/conference\_introduction\_video.html} ) and his overall worldview and perspective on Tradition is similar to ours here at Gornahoor.


\hfill

\texttt{Lapis on 2014-10-27 at 09:24 said: }

I am suprised about Prince Charles. I know that he is a friend of Agha Khan IV whoes grandfather Guenon named as an very important agent of counter-initiation. I listen some of Agha Khans speeches and all i can say is that he is very anti-traditional. And it looks that biography of Prince Charles and his views on many things shows that he is not a man of traditional perspective, he is maybe traditionalist but not traditional and thats a great difference. Maybe they are using Guenon as a weapon for their own needs, just like Alexander Dugin do.


\hfill

\texttt{Tom Blanchard on 2014-10-27 at 18:24 said: }

Lapis,

Which views specifically do you take issue with? From what I've read of his writing, Charles' views on organic farming, of man's need to live in a harmonious relationship with nature, etc. are almost identical to the views expressed by Lord Northbourne in his book, Look to the Land ( \url{http://www.amazon.com/Look-Land-Lord-Northbourne/dp/1597310182} ) Northbourne first translated many of Guenon's books into English, so I think it's fair to say that he also belongs to the same stream of thinking.

Concerning the somewhat tenuous links to ``agents of counter-initiation", Guenon himself affirmed the validity of Freemasonry's initiation, a claim that is rather far from orthodox and might raise similar questions. We can't judge a man by one or two of his friends or by some isolated, seemingly questionable opinion he holds, but rather by the overall thrust and accomplishments of his life and work.


\hfill

\texttt{Lapis on 2014-10-28 at 09:55 said: }

i dont have any problem with views you mentioned, i also support all that stuff but it is insufficient to proclaim somebody a follower of any tradition. Current Dalai Lama also supports all that but he is far from been traditional, a thing that some buddhist monks expirienced on their own skin as he tries to destroy all that is valuable in Buddhism… The Devil is a trickster. As for Freemasonry Guenon stated that only rituals are valid, but that initiation is virtual at the best, and i dont see any problem in it. As for friends of the Prince, i wil say that in my country there is an old proverb which says: ``With whom you are, you are like him". And it is not question of quantity, cause it doesnt matther if there is one or hundred of them. And if he is traditional, than which traditional doctrine is he following? That is what matters, not ``stream of tought". It is not a philosophy. But if you wish to speak about thoughts and his works, then thing that shows him in his real light is his support for modern sport, an i dont want to write why is that so becouse it is obvious for anybody who is fammiliar with the traditional notion of action.


\hfill

\texttt{Cologero on 2014-10-28 at 10:20 said: }

There is no doubt about the Prince's knowledge of, if not commitment to, Tradition.

E.g., \textit{An Introduction from His Royal Highness The Prince of Wales} (video): \url{http://www.sacredweb.com/conference06/conference_introduction_video.html}

and \textit{An Introduction from His Royal Highness The Prince of Wales} (transcript): \url{http://www.sacredweb.com/conference06/conference_introduction.html}

There are also the persistent rumours that the Prince is a convert to Islam. Most likely, he considers the particular traditions not to be of primary importance.

So, yes, probably he follows a ``stream of thought" rather than a specific tradition, but I don't know.

Yet the role of the (future) Monarch has to be taken into consideration. He has to balance different forces within society since he would be the monarch of all his subjects. For that reason alone, he will certainly never prohibit football matches in the UK. And so on.


\hfill

\texttt{August on 2014-10-29 at 08:03 said: }

One is tempted to expect more from men in positions of power, some kind of demonstration of conviction, if they have been touched by real knowledge of the times in more than an intellectual way. But I suppose even a Prince can, like any of us, look around and realise the futility of fighting and simply live his life. 

This highlights another complication, the races of the spirit, differing between individuals, nations and so on. Awkward discussions are avoided in these more academic circles (Sophia journal, Sacred Web etc), with Guenon included because he simply cannot be ignored (though parts of his work can be, as required), while Evola, for example, is excised.

Maybe some of us have acquired such particular tastes with our favourite teachers that everything short of their firmness must necessarily strike us as bland, or compromised.

The time may come when, instead of principles, a clear statement of intent will be the discriminator.


\hfill

\texttt{Synodius on 2014-10-29 at 09:20 said: }

``Maybe some of us have acquired such particular tastes with our favourite teachers that everything short of their firmness must necessarily strike us as bland, or compromised.

The time may come when, instead of principles, a clear statement of intent will be the discriminator."

August, you expressed really well the feeling I always got after some time of reading papers of those `academic circles'. When I asked myself if these are the people who will help to bring about the real restoration I always just left those sites mainly because of boredom. Traditionalism can be used by current `elites' as just another tool and compromised figures as Charles endorsing it are pointing in that direction. Personally I am looking towards Russia (not Dugin) for some lively impulse…


\hfill

\texttt{Cologero on 2015-01-04 at 10:42 said: }

A friend sent me this link re Prince Charles:

The Philosopher Prince: \url{http://www.theamericanconservative.com/articles/philosopher-prince/}

Perhaps, given the expected new encyclical, the Prince will find himself closer to the current pope.


\hfill

\texttt{Tom Blanchard on 2015-01-12 at 12:26 said: }

Cologero: Thanks. They will be even further reconciled when divorce and remarriage are legitimized at the second Synod on the Family session later this year..! (har, har)

Jokes aside, that is indeed quite possible. As difficult as Pope Francis is to stomach (for some of us, at least), his papacy may somehow turn out to be a net benefit for Christianity. The ``super-correct" SSPX type traditionalists miss the point that since the faith has expanded beyond Europe, the answer is not to force non-westerners into European modes of worship, but to establish something along the lines of entirely separate rites for these people (if we need separate rites for the East and West of the Roman empire, how much more so should there be a separate rite for Christianity in the Philippines…). Binding peoples together in a higher unity while acknowledging their earthly diversity, this solution would allow both for the preservation of the purity of traditional European worship as well as for the rest of the world to have forms of worship that are consonant with who they are, how they experience life and spirituality, etc. 

The alternative is the ``lowest common denominator" approach to catholic unity, which results in the sterile, mundane, and superficial form of religion that is imposed almost universally today.


\hfill

\texttt{Cologero on 2015-01-12 at 16:13 said: }

Tom B., those are good points and apropos of something we have planned following the upcoming Evola translation this week. Celsus defended paganism against the Christian accusation of ``polytheism". His defense was that all the local cults were hierarchically arranged with the Platonic God at the top. That seems to be similar to your point. Evola will reprise that same argument. His ineluctable problem, however, is that the local pagan cults are now all spiritually dead.

Curiously, the super-correct would regard your proposal as the LCD approach


\hfill

\texttt{Taivasvaeltaja on 2024-01-27 at 15:41 said: }

``Deity sleeps in a stone, breathes in a plant, dreams in an animal, and comes alive in man."

– Ancient Nordic Proverbs

Preface

A universal esoteric doctrine holds that there is no such thing as inanimate nature. The different kingdoms that exist in nature are all alive in their own way, and in the kingdoms lower than man are influenced by preconscious, collective forces, which are also called mass souls. From the ancient Nordic statement above in the article, we can see what these kingdoms of nature are and how the Deity manifests in them: they are 1) the rock and mineral world, 2) the plant world, 3) the animal world, and 4) humanity. In addition to this, esoteric teaching presents the idea of \hspace{0pt}\hspace{0pt}the so-called fifth kingdom; in other words, we could properly call it true humanity, as the fourth kingdom is the kingdom of the animal-man. Let us briefly examine these five kingdoms of nature one at a time, and their significance in the scheme of planetary evolution.

Stone and Mineral Municipality

Nature's lowest level of development is the stone and mineral kingdom, where the Deity still sleeps completely. The Cosmic Sound, the development impulse of the Logos and world reason manifests itself in this realm motionless and stagnant, according to the initial quote, it is completely asleep. The forces that act in the stone and mineral kingdom are completely unconscious, they are animated only by the various cosmic principles and angelic forces that bring about the richness and diversity of the entire stone and mineral world. The elemental correspondence of the stone and mineral kingdom is Earth and the symbol is Square. In humans, the stone and mineral kingdom manifests itself especially in the bones. The planetary equivalent of the mineral kingdom is Saturn, which, like a skeleton, symbolizes time, aging, and death.

Flora

The next kingdom of nature is the plant kingdom, where the divine development impulse comes to life for the first time and begins to breathe, according to the quote. We can also notice in this expression how, for example, forests are generally called the lungs of the planet earth. The developmental impulse of the Logos takes its first actual living and preconscious form in the flora, where collective, cosmic forces and principles act. In the flora, for the first time, Deity can move. The diversity of nature is also an expression of the action of various cosmic forces. The elemental correspondence of the flora is Water, its symbol is the Chalice or downward pointing triangle. In humans, the plant kingdom manifests itself as a nervous system. The planetary equivalent of the plant kingdom is the Moon.

Animal kingdom

The third kingdom of nature is the animal kingdom, where the divine development impulse, according to the quote, dreams. The development impulse of Logos acts in it again a degree more freely than in the flora. This can be seen as a metaphor for how different instincts and urges begin to work in animals, which guide the animals' actions according to cosmic principles. When we look at an animal and its actions, we easily think that the animal makes its own independent decisions like a human, but in reality the animal is moved because it has no independent reason. The animal kingdom is influenced by different collective souls, which bring about the diversity of the animal kingdom. The animal's elemental correspondence is Fire and the symbol is either a Wand or an upward pointing triangle. In humans, the animal kingdom manifests itself as internal organs and blood. The planetary equivalent of the animal kingdom is Mars, which symbolizes the fire of life.

Mankind

The fourth kingdom of nature is humanity as we know it, where the divine developmental impulse first comes to life as a conscious being. Every human being has individual reason – man is a rational animal – at least to some extent, although a large number of people are still influenced by various animal drives and instincts as well as the collective impulses shared by mankind. One characteristic of humans is that each member of the human race is capable of producing reproductive offspring, whereas various members of the animal kingdom are unable to do so. So it can be said that the striving for unification affects the original dispersion in man, which is also the biological basis of brotherhood and sisterhood; certainly does not mean uncontrolled mixing. The purpose of humanity's level of development is the individualization of each person as far as possible – individuation – where collective effects are even less important. This is the freedom to which Christ sets us free – to transcendent bliss and joy. With its central role, humanity gives its own development impulse to the lower kingdoms of nature; Humanity's actions therefore have a concrete effect on what kind of manifestations appear in the lower kingdoms, e.g. the animal kingdom: for example, predators appear in the world because humanity has not freed itself from atavistic violence. 1 Humanity's elemental correspondence is Air and its symbol is the Sword. 2 Humanity's bodily correspondence is the brain. The planetary equivalent of humanity is Venus-Lucifer.

The Fifth Kingdom

The fifth kingdom is such a level of development that such a person is no longer born into the world, but influences the world from higher levels, inspiring and leading this side of humanity upwards from supersensible levels. The fifth kingdom is the level of development called the Kingdom of Heaven; the people living in this kingdom have previously been part of the fourth kingdom / animal-humanity, and through their own spiritual efforts, after countless incarnations, have reached such a spiritual level in their development that they no longer need to incarnate; they are more human than us humans living in the world, in which many animal traits still live. They have purified themselves completely, ``sacrificed their inner animals" on the altar of spiritual development, freed themselves from their karma and the compulsion of incarnation.3

However, the Master Human can also incarnate if he has a special mission in the world. Even then, he is practically free from his body, which is to him like the Master's own five-windowed House in the Lord's room, where he is free to move as he pleases. A master has a knowledge of the Buddhist field that enables many seemingly supernatural things such as telepathy, telekinesis, diamond creation, alchemy secrets – and so on – depending on his line of development and tradition. The masters guide cultural seasons, e.g. With Manu. According to tradition, there are fourteen Manus in total, and they act as the spiritual centers of each age from the supersensible levels. We can only imagine what kind of tasks the Master can assume in the spiritual hierarchy of our planet and the cosmos. They can also be called Master people – they are fully initiated into the Mysteries, liberated souls. The members of this fifth kingdom have consciously decided to remain connected to the planet earth to help all its kingdoms to develop forward – in particular they inspire the most developed members of humanity operating in the world.4 The elemental correspondence of the Master person is the so-called ether and his symbol is the upward-pointing pentagram, which has long been known as the Master's symbol. The planetary equivalent of the Master Adept is Jupiter.

1 The utopia of the distant golden age, when humanity has been freed from all violence, is expressed in the Bible with the following words: ``Then the wolf will live with the lamb, and the panther will lie down next to the goat; the calf and the young lion and the bullock are together, and the little cub is shepherding them. The cow and the bear go to the pasture, their calves and cubs lie down together, and the red deer eats the fodder like a deer." – Isaiah 11:6-7

2 In the symbol we can see the connection to the principle of reason, which is depicted by the sword.

3 ``In the resurrection, one does not take a wife or become a wife. The resurrected ones are like angels in heaven." – Matt. 22:30

4 This would seem to contradict what I wrote in the text Gone Fishing Beyond the Sea. Let me clarify a bit. From this very real situation, a phenomenon has arisen in the New Age, known as ``channeling of ascended Masters", where a person claims to be channeling the messages of any ``Master" into the world at any time. However, when we look at these channeling phenomena and the content of the messages with the right kind of skepticism, we can very well notice that it is often a question of very confused psychics and practiced by people who have very questionable tendencies towards medium-ship, and the messages concern any matter that is happening in the world at any time, which in reality expresses one's own personal and a subconscious wish for how things would be. How to distinguish true inspiration from psychic medium-ship? I would say that true inspiration is such that a person's own Ego is fully aware and awake, and such a person only so-called ``opens the upper channels" to this spiritual world, from which new inspiration and information flows into him. This kind of inspiration often does not directly concern worldly matters, but instead manifests itself as, for example, scientific breakthroughs or artistic inspiration.


\hfill


\end{sffamily}\end{footnotesize}
